\documentclass[11pt,a4paper]{article}
\usepackage[margin=24mm]{geometry}
\usepackage{amsmath,amssymb,amsthm,hyperref}
\newtheorem{theorem}{Theorem}\newtheorem{lemma}[theorem]{Lemma}
\newtheorem{proposition}[theorem]{Proposition}\newtheorem{corollary}[theorem]{Corollary}
\newcommand{\Sz}{\operatorname{Sz}}\newcommand{\Aut}{\operatorname{Aut}}
\newcommand{\Cor}{\operatorname{Cor}}\newcommand{\gen}[1]{\langle#1\rangle}
\title{Suzuki triality geometries from generating cyclic triples}
\author{Human Compression project: L015 working application}\date{}
\begin{document}\maketitle
This is an unreviewed complete-application draft for the fixed fourteen-page
preprint of Leemans--Stokes--Tranchida~\cite{LST}. We retain their triangle,
chamber-system and coset constructions and their attribution to the classical
Suzuki structure. The accompanying \emph{Generating Suzuki triples by
excluding proper subfields} supplies the uniform generating-seed theorem,
including all its retained matrix and group-structure inputs. The entire
four-page accepted \emph{Coset trialities, chamber recognition and reduced
relations} is a further common dependency; its Theorem 2 supplies the
rank-three recognition argument. Neither common paper is a free import.
No finite calculation or new formalization is a premise of this application.

\section{Local supports and their stabilizers}
Let $G=\Sz(q)$, $q=2^n$, $n>1$ odd, with its ovoid $\Omega$.
In addition to the explicit structure retained in the common paper we use
the classical Zassenhaus property: no nonidentity element fixes three ovoid
points. A nonidentity split-torus element fixes exactly its two endpoints;
a nonidentity nonsplit-torus element fixes none, since a point stabilizer
has odd part dividing $q-1$, coprime to the nonsplit torus orders.
These are the source's classical action inputs, not new classification claims.
For an odd dihedral subgroup $D=\gen{x,y}$ of order $2m$, $m>1$, put
$O(D)=\{\text{fixed point of }r:r\text{ is an involution of }D\}$.

\begin{lemma}[Fixed-point and support recognition]\label{support}
Distinct involutions with the same fixed point commute. Distinct involutions
with different fixed points have a product of nonidentity odd order.
For fixed distinct points $P,X$ and a fixed involution $r$ at $P$, the
products $ru$, as $u$ ranges over the involutions at $X$, have pairwise
distinct conjugacy classes. The support $O(D)$ has $m$ points and determines
$D$ uniquely among dihedral subgroups of order $2m$.
\end{lemma}
\begin{proof}
The involutions at a point are the nonidentity elements of its elementary
abelian root center. Normalize $P,X$ to $\infty,0$, and write
$r=z(b)$, $u=wz(t)w$ with $b,t\ne0$. The common paper's direct trace formula
gives $\operatorname{tr}(ru)=(bt)^\theta$. It is nonzero and varies
injectively with $t$. The retained even-order trace property and invariance
of trace prove both remaining fixed-point assertions. In particular two
different reflections of $D$ cannot share a fixed point, since their
product is a nonidentity odd rotation. The cyclic rotation group $K=C_m$
acts regularly on the $m$ reflections and hence on their fixed points.

Let $A=\operatorname{Stab}_G O(D)$. This is proper: $m$ is at most the largest
torus order, which is less than $q^2+1$. It contains $D$ and cannot lie in a
Borel, because $D$ has involutions with different fixed points. Descend the
maximal-subgroup list. If $A$ equals a subfield Suzuki group $\Sz(r)$ with
$r\ge8$, then $O(D)$ lies in that group's subovoid and is a nonempty invariant
subset. Transitivity makes it the whole subovoid, impossible since
$m\le r+\sqrt{2r}+1<r^2+1$. Thus descent eventually places $A$ in an odd
torus normalizer, with the terminal $\Sz(2)=C_5:C_4$ included in that case.
The rotation subgroup $K$ is in the normalizer's unique odd cyclic subgroup,
so that normalizer, and therefore $A$, normalizes $K$.

Let $T=C_G(K)$ be the ambient cyclic torus. We claim $A\cap T=K$.
For $t\in A\cap T$ and $P\in O(D)$, regularity of $K$ on the support lets
us choose $k\in K$ such that $tk$ fixes $P$. For a nonsplit torus this forces
$tk=1$. For a split torus, $P$ is not an endpoint: a reflection inverting
the torus interchanges its endpoints, since fixing both would put that
involution in the odd two-point stabilizer. The same conclusion follows.
Hence $A/K$ embeds in $N_G(K)/T$, which is cyclic of order two or four.
Its unique order-two subgroup has preimage $D$. Any other dihedral group
of order $2m$ with the same support lies in $A$; its odd rotations lie in
$A\cap T=K$, and it has the same preimage. It therefore equals $D$.
\end{proof}
This proves the intended distinct-input version of source Lemma 2.2; the
printed version must exclude $\zeta=\zeta'$ and require $P\ne Q$.
The first omission gives identical products. For the second, take distinct
$\zeta,\zeta'$ at $P=Q$, both different from the fixed involution $\omega$.
Then $\omega\zeta,\omega\zeta'$ are nonidentity involutions and are
conjugate: point double transitivity and the split-torus scaling show
that all involutions form one class.

Fix an order-three automorphism $\tau$ and an involution $a$, put
$b=\tau(a)$, $c=\tau^2(a)$, and suppose $|ab|=m>1$ is odd. All pair
products have order $m$. Write
\[
 H_0=\gen{b,c},\quad H_1=\gen{c,a},\quad H_2=\gen{a,b},\quad O_i=O(H_i).
\]
The common paper's Sylow-conjugacy argument allows us to normalize $\tau$
to a field automorphism without changing any assertion here.
The triple $(O_0,O_1,O_2)$ is the source triangle: its pairwise intersections
contain the fixed points of $a,b,c$, respectively. Cyclic symmetry implies
that equality of any two supports, or any two $H_i$, implies equality of all.

\begin{lemma}[Nondegeneracy and properness]\label{proper}
The triangle is nondegenerate exactly when the $H_i$ are distinct.
In that case $H_i\cap H_j$ is the common marked involution subgroup.
If $m\equiv2\pmod3$, the triangle is nondegenerate. If $m=5$, each pair
of supports meets in just its marked fixed point, and their triple
intersection is empty.
\end{lemma}
\begin{proof}
The first claim follows from Lemma~\ref{support}. If two different $H_i$
share a nonidentity odd element, its centralizer is a torus containing
both order-$m$ rotation groups. A cyclic torus has only one subgroup of
order $m$; the common marked reflection then makes the two $H_i$ equal.
If they share two distinct involutions, their product is such an odd element.
Thus their intersection is precisely the common order-two subgroup.

If all $H_i$ coincide, the argument in the common paper's descent lemma
gives a unit $u$ modulo $m$ with $1+u+u^2\equiv0\pmod m$. A prime divisor
of $m$ congruent to two modulo three makes this impossible when
$m\equiv2\pmod3$. This proves nondegeneracy without assuming ambient generation.

For $m=5$, suppose $X$ is an additional point common to $O_i,O_j$, besides
the point $P$ of their common marked reflection $r$. Choose the respective
reflections $u,v$ fixing $X$. Both $ru,rv$ have order five. There is a unique
conjugacy class of order five, so the fixed-point trace injection forces
$u=v$. This contradicts the just-proved intersection of the two dihedral
groups. Thus the pair intersection is just $P$. The three marked fixed
points are distinct, since each marked product has odd order greater than
one; the triple intersection is therefore empty.
\end{proof}
The lemma recovers source Lemmas 3.4 and 3.6 without their use of the false
unrestricted generation step. Nondegeneracy alone does \emph{not} imply
$\gen{a,b,c}=G$; that implication in source Corollary 3.5 is separately
refuted by the common paper.

\section{Prime products and the coset geometry}
\begin{lemma}[Prime rotation normalizers]\label{prime-normalizer}
If $K=C_p\le G$ for an odd prime $p$, a dihedral subgroup $D$ of order $2p$
contains $K$ if and only if $D\le N_G(K)$.
\end{lemma}
\begin{proof}
Containment implies normalization because $K$ is characteristic in $D$.
Conversely $N_G(K)$ is its torus normalizer and has a unique subgroup of
order $p$. The rotation subgroup of $D$ must be that subgroup.
\end{proof}

\begin{lemma}[The full four-element product identity]\label{product}
Suppose $m=p$ is prime and the source triangle is proper and nondegenerate.
Then, with every ordered renaming of $a,b,c$ allowed,
\[
 H_1\cap H_2H_0=\{1,a,c,ac\}=\gen{a}\gen{c}. \tag{1}
\]
\end{lemma}
\begin{proof}
The displayed four elements plainly belong to the intersection. We prove
the reverse containment. First take reflections $u\in H_2$, $v\in H_0$.
If they coincide, their product is one. If they are distinct but have the
same fixed point $X$, their product is another nontrivial root involution
at $X$. Membership in $H_1$ would put $X$ in all three supports, contrary
to properness. With different fixed points their product is odd. If it
belongs to $H_1$, it has order $p$ and generates its rotation subgroup $K_1$.
Both $u,v$ invert $uv$, hence belong to $N_G(K_1)$. Also $a,c\in N_G(K_1)$.
If $u\ne a$, the two distinct reflections $u,a$ generate $H_2$ (here
primality is essential), so Lemma~\ref{prime-normalizer} puts $K_1$ in $H_2$.
The common reflection $a$ then implies $H_1=H_2$, a contradiction. Thus
$u=a$; similarly $v=c$. The only reflection--reflection outputs in $H_1$
are consequently $1$ and $ac$.

Next let $u\in H_2$ and $v\in H_0$ be nonidentity rotations. Rewrite
$uv=(ub)(bv)$ as a product of two reflections. The result, if in $H_1$,
is one or $ac$. It cannot be one, since this would give a nontrivial odd
element in $H_2\cap H_0=\gen{b}$. Moreover equality to $ac$ forces
$ub=a$ and $bv=c$, so $u=ab,v=bc$. In particular a product of nontrivial
rotations from two of the $H_i$ cannot be a reflection in the third.
This last assertion holds after every renaming of the indices.

Finally suppose $u\in H_2$ is a nontrivial rotation, $v\in H_0$ a
reflection, and $z=uv\in H_1$. If $z$ is a reflection, the identity
$u=zv$ and the reflection--reflection argument, now with the indices
renamed, force $z=a,v=b$, yielding $u=ab$ and output $a$.
If $z$ is a rotation, $u^{-1}z=v$ contradicts the rotation--rotation
assertion. The reflected case gives output $c$. Products with an identity
factor are covered by the pairwise intersections. These are all elements
of the odd dihedral groups, proving (1).
\end{proof}

\begin{theorem}[Selected generating triangles]\label{geometry}
Under Lemma~\ref{product}, additionally suppose $\gen{a,b,c}=G$.
Then $\Gamma=\Gamma(G;H_0,H_1,H_2)$ is a thin, residually connected,
flag-transitive rank-three geometry with full typed automorphism group $G$.
Every rank-two residue is a $p$-gon.
\end{theorem}
\begin{proof}
Lemma~\ref{proper} gives the three pairwise order-two intersections and
trivial triple intersection. Equation (1) gives the exact product identity
of Theorem 2 in the whole retained common paper on coset trialities.
That theorem proves the geometry, all-flag transitivity, thinness, residual
connectivity and full typed group, with no finite enumeration.
The rank-two residue at $H_i$ is the coset graph of $D_{2p}$ over its two
marked reflection groups, hence a cycle of length $2p$, a $p$-gon.
\end{proof}
Without $\gen{a,b,c}=G$, the same argument applies to the generated subgroup
$H$, and the ambient coset system is a disjoint union indexed by the cosets
of $H$. This is the exact missing generation boundary in source Theorem 3.14
and the final direct-coset paragraph. The subgroup product calculation does
not by itself connect those components.

\section{The literal triangles and chamber-system identification}
\begin{lemma}[The three-point obstruction]\label{three-points}
Let $P,Q,R$ be three distinct ovoid points cycled by $\tau$. Their setwise
stabilizer in $G$ is trivial.
\end{lemma}
\begin{proof}
The pointwise stabilizer is trivial by the Zassenhaus property, so the
setwise stabilizer embeds in $S_3$. It contains no three-cycle because
$3\nmid|G|$. If $g$ induces a transposition, its image under conjugation
by $\tau$ induces a different transposition. Their product is a forbidden
three-cycle. The stabilizer is therefore trivial.
\end{proof}
For a nondegenerate triangle, an element of $G$ stabilizing its unordered
set of three supports permutes the three $H_i$, by Lemma~\ref{support},
and therefore permutes their pairwise marked involutions, by
Lemma~\ref{proper}. It stabilizes their three fixed points and is the
identity by Lemma~\ref{three-points}. Thus both the ordered and unordered
triangle orbits have size $|G|$, proving the full source Corollary 3.10.
The distinct-cycle interpretation is essential in source Lemma 3.8.

Use left action notation temporarily: $g$ carries the base triangle $T$
to $gT$, and carries a subgroup $H_i$ to $gH_i g^{-1}$. The fixed-point
sets transform in exactly the same way, proving source Lemma 3.9 with
its conjugated pair $(gag^{-1},g\tau g^{-1})$. The unique reflection
attached to type $i$ at $gT$ is $gr_i g^{-1}$, where $(r_0,r_1,r_2)=(a,b,c)$.
The source adjacency rule therefore joins $gT$ to
$(gr_i g^{-1})(gT)=gr_iT$. It is exactly right multiplication by $r_i$
on the chamber labels $G$. This proves symmetry, the two-element panels,
equivariance and source Theorem 3.12 without a hidden connectedness premise.
Its components are $g\gen{a,b,c}$.

The type-$i$ objects of the associated chamber incidence system are the
components under the other two colors, namely $gH_i$. Two such objects
are incident exactly when these subsets meet. Inversion sends these left
cosets to the right cosets $H_i g^{-1}$ of Theorem~\ref{geometry}; it
preserves nonempty intersections. This is an explicit identification
with the source geometry, rather than only an equality of diagrams or groups.

\section{Full correlations and hypermaps}
\begin{proposition}[Normalized correlations]\label{correlations}
Let a thin geometry have full typed group $G$ acting regularly on chambers.
Fix a chamber $C$, and write $r_i$ for the unique group element sending
$C$ to its other type-$i$ panel chamber. Every correlation with type
permutation $\pi$ induces an automorphism of $G$ sending $r_i$ to $r_{\pi(i)}$.
If its chamber graph is connected, the normalized chamber stabilizer
embeds in the permutation group on the types.
\end{proposition}
\begin{proof}
Multiply the correlation by the unique typed automorphism that returns
its image of $C$ to $C$. Call the result $\beta$. The typed group is the
kernel of the type action, hence is normalized by $\beta$. Conjugation
by $\beta$ is thus an automorphism of $G$. The element
$\beta r_i\beta^{-1}$ sends $C$ to its unique other type-$\pi(i)$ panel
chamber, so is $r_{\pi(i)}$. This proves source Proposition 3.15 without
assuming the chamber graph connected or asserting that these $r_i$
already generate $G$. In the connected case, a typed automorphism fixing
$C$ fixes its unique neighbors and then every chamber along all paths;
this proves the claimed embedding. The source's preceding assertion
that all $r_i$ generate $G$ needs this connectedness hypothesis; the
proposition itself does not.
\end{proof}

\begin{theorem}[The complete main geometry conclusion]\label{main}
For every odd $n>1$ divisible by three, $G=\Sz(2^n)$ is the full typed
automorphism group of a thin, residually connected, flag-transitive rank-three
geometry with each of the three rank-two residues a pentagon. Its full
correlation group is $G\rtimes C_3$, with trialities and no dualities.
\end{theorem}
\begin{proof}
Choose the full generating seed of order five from the uniform common
theorem. Lemma~\ref{proper} makes its triangle proper and nondegenerate.
Theorem~\ref{geometry} supplies every claimed geometric property and
the full typed group. The actual order-three field automorphism sends
$H_i$ to $H_{i+1}$; on right cosets it acts as
$H_i g\mapsto H_{i+1}\tau(g)$. The first proposition of the common
coset paper proves this is a correlation of order exactly three.

Normalize any correlation to fix the standard chamber. Its type action
is faithful, by Proposition~\ref{correlations}. If a transposition occurs,
the corresponding normalized correlation $\beta$ has square the identity.
Conjugation on $G$ is then an automorphism of order dividing two and
interchanges two of the distinct marked involutions. Since
$\operatorname{Out}(G)=C_n$ has odd order, this automorphism is inner.
Its conjugating element permutes the three marked fixed points, so is
the identity by Lemma~\ref{three-points}, a contradiction. The type image
contains $C_3$ and no transposition, hence equals $C_3$. The displayed
field automorphism splits the extension and proves $\Cor\Gamma=G\rtimes C_3$.
This is the existential source Theorem 3.16, with its generation proof repaired.
\end{proof}

More generally, for every odd element order $m\equiv2\pmod3$ the uniform
common theorem gives a regular hypermap of type $(m,m,m)$ with group $G$:
let the three free-product involutions act on the flag set $G$ by right
multiplication by $a,b,c$. Generation makes this action transitive, and
its full commuting automorphism group is left multiplication by $G$.
The field automorphism gives the triality. A type-permuting symmetry can
be normalized to fix the identity flag and then acts on all words in the
three generators; it induces the same marked group automorphism as above.
The three-point argument therefore again excludes a duality, and the full
type-permuting symmetry group is $G\rtimes C_3$. No unproved flag-transitivity
of a general-$m$ coset geometry is used for this hypermap conclusion.
Exactly one of $q+\sqrt{2q}+1$ and $q-\sqrt{2q}+1$ is two modulo three:
$q\equiv-1\pmod3$ and $\sqrt{2q}$ is $1$ or $-1$ modulo three.
This recovers the source's full-torus hypermap choices and its stronger
existential assertion for every realized $m\equiv2\pmod3$.

\section{Corrections and retained background boundaries}
Apply the written uniform theorem in $K=\Sz(8)$ with inverse field
automorphism and $m=5$, and embed it as the standard subfield of
$G=\Sz(2^{15})$. The ambient automorphism $F^5$ restricts to $F^2$ on $K$.
The triple therefore generates exactly $K$, not $G$, and has distinct
rotation groups and a moved marked fixed point. Lemma~\ref{proper} proves
that its triangle is proper and nondegenerate. The source triangle orbit
still has $|G|$ chambers, but its chamber graph has $[G:K]>1$ components.
This is an exact written counterexample to the unrestricted second
assertion of Theorem 3.2, Corollary 3.5, the ``any triangle'' ambient-group
part of Corollary 3.7, and Theorem 3.14 as printed. It also qualifies the
final claim that ambient residual connectivity is straightforward.
The positive replacements above require a selected fully generating seed;
they do not identify the source's literal theorem set with a true theorem set.

For the source's geometric warning, $\operatorname{Stab}_G O(D_{10})$
has order twenty, not ten. Indeed five divides a nonsplit torus order.
In its normalizer $T:C_4$ every involution is the square of an order-four
element, because it lies in a Sylow $2$-subgroup, which is cyclic of order four.
Given the reflection $r$ of $D_{10}$, choose such a square root $x$.
The element $x$ normalizes its unique $C_5$ and centralizes $r=x^2$,
so normalizes $D_{10}$ and preserves its support. Lemma~\ref{support}
now gives the full stabilizer $C_5:C_4$, of order twenty, isomorphic
to $\Sz(2)$. Thus the source's reason for retaining chamber information
rather than using the bare five-point subsets is exact.

The general chamber representation and connectivity facts in source
Proposition 2.3 and Lemma 2.4 have elementary orbit proofs: send $gB$ to
$gC$ and use its stabilizer $B$; panels are the corresponding $G(i)$-cosets;
a word in the panel subgroups traces a gallery and every gallery gives
such a word. Source Theorems 2.5 and 2.6 are explicitly retained general
chamber/geometry correspondence theorems from Buekenhout--Cohen~\cite{BC},
not silently proved by our rank-three criterion. In the former, for a
nontrivial Borel use $G(J)=\gen{B,G(j):j\in J}$, in particular $G(\varnothing)=B$;
the printed empty-generated-subgroup convention otherwise loses $B$.
For example, $G=B=C_2$ with both panel subgroups equal to $G$ gives a
connected one-chamber, residually connected system. The printed identity
with $L=\{0\}$, $J=\{1\}$, $K=\varnothing$ would say $G=1$.
Our application has $B=1$ and directly proves the required rank-three
correspondence, so no unrestricted presentation of that background theorem
is needed as an additional application premise. A coset \emph{pregeometry}
is not called a geometry before the flag-extension condition is proved.
The historical census numbers and earlier-priority remarks in the introduction
remain attributed literature reports, not new enumerations or prerequisites.

\begin{thebibliography}{9}
\bibitem{LST} D. Leemans, K. Stokes and P. Tranchida,
\emph{Flag transitive geometries with trialities and no dualities coming
from Suzuki groups}, arXiv:2311.13522v1, 22 November 2023, fourteen-page
fixed edition. The source construction and all source-result attributions.
\bibitem{BC} F. Buekenhout and A. M. Cohen, \emph{Diagram Geometry},
Springer, 2013, Propositions/Theorems 3.6.4, 3.6.7, 3.6.9 and 3.4.6,
as cited in the fixed source. General background retained with the
explicit Borel convention; no local complete external proof charge is known.
\end{thebibliography}
\end{document}
